% !PRINTED 695
\weberpage{695}
%<<CONT
wird, und wenn man hier $m=p$ setzt, so folgt, daß es Klassen
der Ordnung $p$ von der zweiten Gattung gibt.

\begin{center}
§ 199. \gesperrt{Der Riemann-Rochsche Satz.}
\end{center}

Das unter dem Namen des \gesperrt{Riemann-Rochschen} Satzes
bekannte Theorem hat den Zweck, die Dimensionen von Polygonklassen
 zu bestimmen.

Wir betrachten zunächst eine \gesperrt{eigentliche Polygonklasse}
 $A$ von der Ordnung $n$, nehmen darin zwei teilerfremde
Polygone $\mathfrak{A}$, $\mathfrak{A}'$ und setzen:
\[
z = \frac{\mathfrak{A}'}{\mathfrak{A}}.
\]

Ist dann $\mathfrak{A}''$ ein anderes Polygon derselben Klasse, so
setzen wir
\[
\omega = \frac{\mathfrak{A}''}{\mathfrak{A}},\quad \frac{\omega}{z} = \frac{\mathfrak{A}''}{\mathfrak{A}'}.
\]

Es ist also $\omega$ eine ganze Funktion von $z$, und $\omega:z$ eine
ganze Funktion von $1:z$. Wenn also $\omega$ nicht konstant ist, so ist
der Exponent $r$ von $\omega$ gleich 1. Umgekehrt ist auch jede ganze
Funktion $\omega$ von $z$ vom Exponenten 0 oder 1 in der Form
$\mathfrak{A}'':\mathfrak{A}$ darstellbar.

Wir erhalten also die ganze Klasse $A$ aus der Gesamtheit
der ganzen Funktionen von $z$, deren Exponent $\leqq 1$ ist. Diese
erhalten wir aber leicht aus der Normalbasis $(\lambda_1, \lambda_2, \ldots, \lambda_n)$ in $z$,
mit der Reihe der Exponenten $r_1, r_2, \ldots, r_n$ (§ 193). Ist $\lambda_s$ die
letzte dieser Funktionen, deren Exponent gleich 1 oder 0 ist, so
ist jede ganze Funktion, deren Exponent $\leqq 1$ ist, in der Form

\weq{1}{\omega = c_0 z + c_1\lambda_1 + c_2\lambda_2 + \cdots + c_s\lambda_s}

darstellbar ($c_0$ muß konstant sein, weil es nach der Definition
der Normalbasis eine ganze Funktion sein muß, und wenn sein
Exponent positiv wäre, so wäre der Exponent von $\omega$ größer als
eins). Demnach ist

\weq{2}{(O, A) = s + 1,}

und diese Zahl ist also immer $\leqq n + 1$.

Die obere Grenze $n+1$ wird nur dann erreicht, wenn alle
$r_2, r_3, \ldots, r_n = 1$ und folglich $p = 0$ ist.

Da in einer eigentlichen Klasse mindestens zwei Polygone
enthalten sind, also $(O, A) \geqq 2$ sein muß, so kann $n = 1$ nur
in dem Falle $p = 0$ vorkommen.
